% $Id: framework.tex 1784 2012-04-27 23:29:31Z  $
% !TEX root = main.tex

\chapter{Framework Architecture}\label{chap:framework}

% Figure placeholder: Framework architecture diagram
% TODO: Add TikZ diagram showing the four-layer architecture
% (Theory -> Payload -> Backbone -> CRDT Implementations)

\section{Overview of the Layered Architecture}

The verification framework is organized as a stack of four logical layers: the Theory layer at the foundation, the Payload layer above it, the Backbone layer next, and concrete CRDT implementations at the top. Each layer builds upon the abstractions of the layer below it, maintaining a clear separation of concerns. At the base is the \emph{Theory layer}, which provides the mathematical foundation (join-semilattice laws and abstract CRDT properties). Above it, the \emph{Payload layer} realises generic state-based CRDT semantics (local state and merge operations). Next, the \emph{Backbone layer} adds replication mechanics, including vector clocks and causal message delivery, to support operation-based CRDTs. Finally, the top layer consists of \emph{concrete CRDT implementations}, which plug into the framework via lightweight adapters. This layered design isolates different aspects of the system: algebraic correctness, state semantics, replication protocol and application-specific logic are all handled in separate modules. Such modularity makes the framework extensible and eases reasoning, since each layer can be specified and verified in isolation and then composed to obtain full end-to-end guarantees.

The layers interact in a disciplined way. The theory layer specifies the convergence requirements that any CRDT must satisfy; the payload and backbone layers implement a generic runtime that will enforce these requirements, provided the CRDT’s state and update functions meet the theory’s axioms. Concrete CRDTs are integrated last, by instantiating the abstract definitions with their own state sets and operations. Because of the strict separation, changes in one layer (for example, using a different network protocol in the backbone, or adding a new CRDT type) do not require re-proving properties of the other layers. In the remainder of this chapter, we detail each layer of the architecture. Section~\ref{sec:theory-layer} covers the foundational theory module, Section~\ref{sec:framework-payload} describes the payload semantics, Section~\ref{sec:framework-backbone} discusses the backbone for causal consistency, and Section~\ref{sec:crdt-adapters} explains how concrete CRDT implementations are connected via adapters.

\section{Theory Layer: \texttt{SemilatticeTheory} and Abstract CRDT}
\label{sec:theory-layer}

At the foundation of the framework is a formal semilattice theory that captures the algebraic conditions for convergence. This theory layer (developed in Chapter~\ref{chap:theory}) defines an abstract CRDT specification through a set of five axioms. These axioms are: (i) \emph{commutativity}, (ii) \emph{associativity}, and (iii) \emph{idempotence} of the merge operator (the three properties making the merge a join-semilattice operation), (iv) existence of an \emph{identity element} (an initial state that acts as a neutral element for merge), and (v) \emph{monotonicity of updates} (every update must only increase the state according to the semilattice order). When a data type’s state and update functions satisfy these properties, it guarantees that all replicas will reach the same state as soon as they have delivered the same set of updates, regardless of order. In other words, the axioms imply the Strong Eventual Consistency (SEC) property of CRDTs, as established by Shapiro \emph{et~al.}~\cite{shapiro2011conflict}. The role of this layer is critical: it provides a reusable mathematical foundation that all CRDT instances share. Key convergence theorems (for example, that merge operations commute to a common outcome) are proven once at this abstract level. Later in Chapter~\ref{chap:implementations}, we will see that each concrete CRDT can inherit these theorems by showing that its particular state set and operations instantiate the Theory layer’s axioms. This approach follows a \emph{theory/adapter pattern}: the heavy lifting is done in the general theory, and each CRDT needs only to provide a small proof (via an adapter) that the theory’s conditions hold in its case. By separating the mathematical model from any operational details, the Theory layer enables a high degree of proof reuse and confidence: as soon as a new CRDT is proven to form a semilattice with monotonic updates, all the convergence results from Chapter~\ref{chap:theory} immediately apply to it.

It is worth noting that the foundation also includes some general-purpose utilities that support the upper layers. For example, the framework provides library modules for finite maps (to represent mappings from replica identifiers to states or timestamps), multisets (to represent network message buffers), and a formalisation of vector clocks (finitely supported maps from replica IDs to counters). These components, while not part of the core CRDT theory, form a useful substrate for building the payload and backbone layers. In particular, a vector clock structure (introduced in Chapter~\ref{chap:vectors} as an example of a join-semilattice on timestamp vectors) will be used in the Backbone layer to track causal dependencies. With the Theory layer’s axioms and these utility structures in place, we can build the executable layers of the framework confident that their design adheres to proven convergence criteria.

\section{Payload Layer: State-Based Semantics}
\label{sec:framework-payload}

The Payload layer implements the core semantics of state-based CRDTs on top of the abstract theory. This layer is realised by the \texttt{PayloadFramework} module (detailed in Chapter~\ref{chap:payload}), which defines how CRDT states are maintained and how they evolve through message passing. Importantly, the payload framework remains parametric over the actual CRDT state and update types: it assumes nothing about the concrete data structure beyond the semilattice properties ensured by the Theory layer. Instead, it provides a generic execution model that any semilattice-based CRDT can plug into. The key components of the payload layer are as follows.

\paragraph{System configuration.}
A global state of the system is represented by a configuration record containing (a) a set of per-replica states and (b) a multiset of undelivered messages (the network buffer). Each replica’s state lies in the semilattice $S$ of the CRDT, and a message contains a piece of state (often an incremental update or a full state transfer) destined for some replica.

\paragraph{Update and merge operations.}
The CRDT’s abstract update function (often called $\mathsf{apply}$) and merge operator ($\sqcup$) are incorporated via an interface called \texttt{PayloadOps} (the concrete payload operations). The payload layer uses these operations to manipulate the states. Notably, the merge operation used here is exactly the semilattice join $\sqcup$ from the Theory layer. Because $\sqcup$ is commutative, associative and idempotent, merges of states can be done in any order and yield a monotonically non-decreasing state. The monotonic update axiom from the theory ensures that any $\mathsf{apply}$ operation moves a state upwards in the lattice (no update can undo information).

\paragraph{Message buffering (enqueue).}
The payload layer provides an \texttt{enqueue} function to model message sending. When a replica wants to send its state (or an update) to another replica, an entry is inserted into the network buffer (message multiset). This does not immediately affect any replica’s state; it simply makes the message available for later delivery.

\paragraph{Message delivery (reduce).}
The core transition in the payload semantics is the \texttt{reduce} step, which models delivering one pending message to its destination. If the message buffer is non-empty, a message is removed and its payload is merged into the destination replica’s state. In code, the destination’s current state $s$ is looked up, and then updated to $s' := s \sqcup \mathit{payload}$ (using the CRDT’s provided merge operation or, equivalently, using an abstract \texttt{payload-apply} function that applies the payload as an update). The updated state is stored back into the configuration, and the message is removed from the buffer. If no message is available, \texttt{reduce} is a no-op (the configuration remains unchanged). Repeatedly reducing messages in an arbitrary order will eventually deliver all messages and merge all updates into every replica’s state.

The payload layer thus defines an execution model for state-based CRDTs: replicas independently apply local updates (which are monotonic) and occasionally exchange states by enqueuing and reducing messages. Thanks to the semilattice properties, the order in which messages are delivered does not affect the final outcome: any sequence of \texttt{enqueue} and \texttt{reduce} steps that delivers the same set of messages will result in all replicas converging to the same state. This is exactly the Strong Eventual Consistency guarantee formalised in the Theory layer. The payload framework’s design ensures that the algebraic assumptions (ACI laws and monotonicity) are respected at runtime: only merges are used to combine states, and states never decrease. Any CRDT-specific logic (for example, what constitutes an update or how to interpret a message’s payload) is abstracted behind the \texttt{PayloadOps} interface. In Chapter~\ref{chap:payload}, we provide the full formal definition of the \texttt{PayloadFramework} module, including its axioms and inference rules, and prove that under the given assumptions it guarantees convergence.

Overall, the Payload layer can be seen as an implementation of a traditional state-based CRDT model (sometimes called a \emph{state transfer} model): each message contains a state to be merged, and the framework takes care of buffering and delivering these state messages. The data flow in this layer, from the abstract algebraic model down to the runtime message processing, is detailed in Chapter~\ref{chap:payload}. It illustrates how a mathematical semilattice $(S,\sqcup,\leq)$ with a monotonic $\mathsf{apply}$ function is instantiated as a concrete payload module (\texttt{PayloadOps}), which in turn is plugged into the configuration and message handlers (\texttt{enqueue}/\texttt{reduce}) to yield an executable CRDT. This flow of information, from the algebraic specification to the operational semantics, is what enables us to connect high-level convergence proofs to low-level behaviour. The payload layer, in summary, provides a generic, verified engine for executing CRDT state merges; it assumes the axioms from the Theory layer, and on that basis it proves fundamental properties like convergence (any two replicas that receive the same updates, in any order, reach an equivalent state). Concrete CRDTs that do not involve additional metadata (for instance, a Grow-Only Counter or Set) can be fully realised at this layer alone.

\section{Backbone Layer: Causal Metadata and Replication}
\label{sec:framework-backbone}

While the payload layer is sufficient for purely state-based CRDTs, many CRDT designs (especially operation-based ones) require tracking causal ordering of updates to ensure consistency. The Backbone layer extends the framework with a replication protocol that incorporates causal metadata. This layer is implemented by the \texttt{CRDTBackbone} module (described in Chapter~\ref{chap:backbone}), which builds on the payload semantics but adds a vector clock service and modified messaging to enforce causal delivery of operations.

In the backbone layer, a system configuration is enriched with a vector clock registry in addition to the state map and message buffer. For each replica (identified by some $\mathsf{Ide}$), the registry stores a vector timestamp summarising the updates that replica has seen so far. Concretely, if $VC$ denotes the domain of vector clocks (for example, functions from replica IDs to $\mathbb{N}$, with all entries initially $0$), then the configuration now consists of three components:
\[
  (\mathit{states}: \mathsf{Ide}\to S,\;
   \mathit{clocks}: \mathsf{Ide}\to VC,\;
   \mathit{net}: \{\text{messages}\}).
\]
Each message in the network is also augmented with a timestamp: a message is now a triple $(i, \mathit{payload}, clk)$ indicating that it is an update being sent to replica $i$ with an attached vector clock $clk$ (the sender’s state of knowledge at send time).

The backbone’s operational semantics include all the functions of the payload layer (since it extends that layer) and add a few more.

\paragraph{Clock tick and broadcast.}
Before a replica broadcasts an update, it performs a \emph{tick} on its own clock. The function \texttt{cfg-tick} increments the local replica’s entry in the clock registry. This models the generation of a new update event (assigning it a unique timestamp). After ticking, the replica’s update can be applied to its state (using the payload’s $\mathsf{apply}$) and then enqueued as a message.

\paragraph{Timestamped messages.}
Sending a message now captures the sender’s timestamp. When a replica enqueues a message in the backbone layer, it includes the current vector clock value from the sender’s registry as the $clk$ in the message. This effectively tags the update with the causality context in which it was generated.

\paragraph{Causal delivery condition.}
The backbone layer modifies the \texttt{reduce} (message delivery) operation to respect causality. A message can only be delivered to replica $i$ when its attached timestamp $clk$ is causally ready, meaning $clk$ does not reflect any update that $i$ has not yet seen. In practice, this is checked by comparing $clk$ to replica $i$’s current vector clock. The condition is that for every component except possibly $i$’s own, $clk[j] \le \mathit{clocks}_i[j]$; in other words, the incoming message’s timestamp is not ahead of $i$ in any other replica’s position. If the condition is not met, the message must wait in the buffer until causally preceding messages have been delivered. This guard prevents violations of causality (deliveries that would make it appear that an update happened without its cause).

\paragraph{Causal message delivery.}
When a message is deliverable, the backbone \texttt{reduce} operation will deliver it, similar to the payload \texttt{reduce} but with two merges: it merges the incoming payload into the state and merges the incoming timestamp into the local clock. Specifically, if a message $(i, \mathit{payload}, clk)$ is being delivered to replica $i$, the framework looks up $i$’s current state $s$ and current clock $vc$. It computes the new state $s' = s \sqcup \mathit{payload}$ (using the payload’s merge, just like before) and also computes a new vector clock $vc' = vc \sqcup clk$ where the join on clocks is taken componentwise (essentially, $vc'[j] = \max(vc[j], clk[j])$ for each replica $j$). The state map is updated with $s'$ for replica $i$, the clock registry is updated with $vc'$ for $i$, and the message is removed from the network. Merging the clocks ensures that replica $i$ records that it has now seen all events that were reflected in $clk$. Because the vector clock join is idempotent and commutative (pointwise maximum), applying the same message’s clock multiple times or in different orders has no ill effect: each update to the clock registry simply brings it to the greatest values seen so far for each component.

With these mechanisms, the backbone layer implements a causal broadcast protocol on top of the payload CRDT semantics. Every update is timestamped with a vector clock, and every delivery is conditional on causal readiness. This guarantees that if one update happened-before another (in the sense of distributed events), then all replicas will apply them in that causal order. However, updates that are concurrent (no causality between them) may be delivered in any order, and thanks to the commutativity of the payload merge, the states will still converge. The backbone ensures exactly-once, causally ordered delivery of commutative updates, which is precisely the model required by operation-based CRDTs. In effect, the Backbone layer glues a causality management mechanism (vector clocks and delivery guards) onto the Payload layer’s order-insensitive merge semantics.

It is important to highlight how the Backbone layer leverages the Theory layer as well. The vector clocks used here form a semilattice in their own right: given two vector timestamps, the componentwise maximum yields their join, and this operation is commutative, associative and idempotent. In fact, the set $VC$ of vector clocks (of appropriate dimension) ordered by $\leq$ (defined componentwise) is a distributive lattice, as discussed in Chapter~\ref{chap:vectors}. The backbone combines this clock lattice with the payload’s state lattice. We effectively work in the product lattice $S \times VC$: an element of the product is a pair $(s, v)$ of a payload state and a vector clock. The join of two such pairs is defined as
\[
  (s_1, v_1)\sqcup(s_2, v_2) = (\,s_1\sqcup_S s_2,\; v_1\sqcup_{VC} v_2\,),
\]
that is, the payload states are joined by their $\sqcup$ and the clocks by pointwise maximum. This product construction is again a join-semilattice (it inherits ACI properties from each component). As a result, the entire backbone configuration space $S \times VC$ still satisfies the convergence axioms. In other words, despite the added complexity of tracking clocks, the backbone layer’s operations preserve the fundamental algebraic structure needed for SEC. This modular composition is what allows us to prove convergence in a piecewise way: Chapter~\ref{chap:backbone} shows that the Backbone layer correctly maintains causal delivery, and because the combined state+clock forms a semilattice, the convergence theorems from Chapter~\ref{chap:theory} still apply to the system as a whole. We thereby obtain a causal CRDT framework where any payload that was SEC in the state-based sense remains SEC even in an operation-based setting with concurrency and asynchrony.

In summary, the Backbone layer adds the necessary machinery to turn the state-based payload semantics into a full replicated system with messages, while upholding order requirements for operation-based CRDTs. It introduces minimal additional assumptions (primarily, that message channels eventually deliver messages and that clocks are synchronised via max-hold semantics) and no new algebraic laws beyond those already in the semilattice theory. The next section discusses how concrete CRDTs make use of these layers via adapters.

\section{CRDT Adapters and Instantiation}
\label{sec:crdt-adapters}

The final layer of the architecture is where concrete CRDT implementations come into play. Rather than building each CRDT from scratch, the framework uses an adapter pattern to instantiate the generic payload/backbone machinery with the details of a particular CRDT. A CRDT adapter is a thin module that connects a CRDT’s native state and operations to the interfaces expected by the Payload (and optionally Backbone) framework.

To create a new CRDT within our framework, one follows a systematic process:

\begin{enumerate}
  \item Define the state and update types for the CRDT. For example, for a Grow-Only Counter (GCounter), the state might be a natural number (or a map of replica IDs to numbers), and the update is an increment operation.
  \item Provide the merge and apply functions. In the GCounter, the merge operation is typically the pointwise maximum of two counters (since each replica’s count is taken as the maximum of what the replicas have seen). The update (apply) function for an increment is just addition on the local replica’s component. Similarly, for a Grow-Only Set (GSet), the merge is set union and the update is an element insertion. Each CRDT’s design dictates these two key operations.
  \item Prove that the five framework axioms hold. The adapter must ensure that the provided merge is commutative, associative, idempotent, has an identity element and that the update is monotonic with respect to the merge. In practice, these proofs are straightforward because they boil down to simple algebraic laws of the chosen data type. For instance, for GCounter we reduce all required properties to facts about $\max$ over $\mathbb{N}$. We know that for natural numbers, $\max$ is commutative, associative, idempotent and has identity $0$ (zero is the least element). Moreover, incrementing a counter is monotonic (if $n$ is a state, then $n \leq n+1$ in the induced lattice order). Thus GCounter fits the framework’s theory perfectly. These simple laws (often referred to as ACI + identity + monotonicity) are usually either obvious or well-documented properties of the data structure (for example, set union has those same properties with identity the empty set). The heavy lifting has already been done in the Theory layer: the adapter’s job is only to verify that the CRDT at hand conforms to that theory.
  \item Instantiate the framework modules. Once the CRDT’s operations are validated against the axioms, we create a concrete instance of \texttt{PayloadOps} (and if using the backbone, an instance of the combined backbone payload operations that includes clock management). This involves writing a small adapter module that defines the CRDT’s specific \texttt{merge}, \texttt{apply}, initial state and so on, and then uses the generic \texttt{PayloadFramework} and \texttt{CRDTBackbone} to obtain the executable specification. Essentially, the adapter tells the framework “treat my CRDT’s merge as the framework’s merge, my update as the framework’s apply,” and so forth. Because the adapter provides concrete functions in place of the abstract ones, we often call this step an instantiation of the framework for that CRDT.
\end{enumerate}

This adapter approach has several benefits. First, it dramatically reduces verification effort. All the core theorems about convergence, idempotence of merges and similar properties have already been proven in the abstract framework. A new CRDT does not require re-proving those from scratch; one only checks the axioms for that CRDT’s specific operations. For example, once we have proven a general convergence theorem in Chapter~\ref{chap:theory}, the GCounter inherits that theorem immediately after we show that $(\mathbb{N}, \max)$ is a join-semilattice and increment is monotonic. The proofs for each adapter are typically concise, as they often reuse library lemmas (for instance, a lemma that $\max$ on naturals is commutative, or that set union is idempotent). In our implementation, we provide theory instantiation files (for example, \texttt{gcounter\_payload\_theory.ath}) that discharge the five axioms for each CRDT’s operations. Once those checks are done, the framework guarantees the CRDT will behave correctly in a replicated setting.

Second, the adapter pattern enforces a consistent structure across all CRDT implementations. Every CRDT in our framework—counters, sets, registers and so on—has a similar shape of specification: a state type, a merge (join) operation, an update function and an initial state. This makes it easy to add new CRDTs by analogy. For instance, implementing a Positive-Negative Counter (PN-Counter) can be done by composing two GCounter adapters (one for increments, one for decrements) and showing that the pair of naturals with pairwise maximum forms a valid semilattice (which it does, as a direct product of two $\max$-semilattices). The framework naturally supports such product constructions due to the composition properties discussed earlier: if the components satisfy the semilattice axioms, so does their combination. In fact, complex CRDTs like the Two-Phase Set (which can be seen as a combination of a grow-only set and a grow-only tombstone set) are handled by providing adapters for each part and then combining them, with minimal new proof obligations. This systematic reusability is a direct result of the layered design.

Finally, the adapters serve as a bridge from abstract theory to executable code. They connect the mathematical model to concrete operations that can run. Because our framework is built in the Athena proof assistant, these specifications are not only proven correct but can also be executed (for simulation or testing). The same adapter that we prove correct is used to generate a running instance of the CRDT in a demo or test harness. This tight feedback loop helps ensure that the verified model truly corresponds to an implementable artefact.

In Chapter~\ref{chap:implementations}, we showcase six verified CRDTs built with this framework (GCounter, GSet, PN-Counter, Two-Phase Set, OR-Set and LWW-Register). Each is defined by a small adapter module that plugs into the Payload layer (for state-based CRDTs) or the combined Payload+Backbone layers (for those that need causal delivery). The result is that each of these CRDTs comes with strong convergence guarantees: since their adapters satisfy the Theory layer axioms, the Payload/Backbone execution will ensure SEC and other correctness properties proven in the base framework. The adapters encapsulate all the CRDT-specific logic, while the heavy verification burden is carried by the general framework. This demonstrates the power of the theory/adapter approach in achieving modular verification: we verify generic convergence once, and then validate each new CRDT via a brief conformance proof. As a consequence, adding a new CRDT to the system is roughly linear in the size of the CRDT’s specification (instead of requiring quadratic or worse effort to develop a full proof from scratch for every new type). The layered architecture, with clear interfaces between Theory, Payload, Backbone and CRDT adapters, thus achieves the thesis goal of reusability and scalability in verifying CRDT convergence.

\endinput
